\section{A finite conical two-colour solver, with its exact curl hypotheses}
\label{sec:r11-conical-euler}

\paragraph{Status and scope.}
\textbf{PROVED THIS ROUND}, submitted for exact-source nonconstructor
review.  The construction contributors are the R11 root and cube agents.
The principal assertion is a finite algebraic solver in every carrier
codimension.  It requires two displayed pure-colour curl identities;
it does not deduce them in higher rank from closedness.  No contraction
on arbitrary Berkovich actual forms is introduced.  The formulas are
first verified in a specified native conical presentation and their
outputs are finite expressions in actual forms.

\subsection{The finite native conical algebra}

Work on a connected conical open $\Omega\subset\mathbf R^d$ with a
fixed finite central polyhedral refinement.  The application permits
$\Omega=\{\tau>0\}$ with all remaining coordinates unrestricted.
Supply finitely many continuous homogeneous piecewise-linear functions
$h_1,\ldots,h_t$, linear on that refinement, and form the native
actual algebra generated by
\[
 h_i,\qquad d'h_i,\qquad d''h_i,\qquad Dh_i=d'd''h_i.
 \tag{CE1}
\]
Finite polynomial expressions and finite common refinements are meant.
There is no completion or infinite sum.  A scalar function, its first
derivatives, and its corner form have their usual native actual
presentations.  Additional homogeneous linear coordinate functions
may be adjoined to the same finite list.

Let $E=\sum_jx_j\partial_{x_j}$ be the radial vector field.  On
these expressions define odd contractions by the finite rules
\[
\begin{array}{c|cccc}
 &h_i&d'h_i&d''h_i&Dh_i\\\hline
 I'&0&h_i&0&0\\
 I''&0&0&h_i&0.
\end{array}
\tag{CE2}
\]
Extend them by the total-graded Leibniz rule.  All final uses of a
contraction mean evaluation by (CE2); thus the result is again a
finite expression in (CE1), not a new type of object.

\begin{lemma}[The finite contraction respects native actual relations]
The operations (CE2) are well-defined on the supplied native actual
algebra.  They agree on every common central carrier with contraction
of its tangential coefficient by $E$.  In particular they preserve
every equality of the supplied native actual expressions.
\end{lemma}

\begin{proof}
On a cone where $h_i$ is linear, $d'h_i(E)=d''h_i(E)=h_i$ by
homogeneity.  Every pure corner carrier of $Dh_i$ is central, so
$E(x)$ belongs to its tangent space at every point of that carrier.
Its normal prime--double-prime current is therefore killed by either
contraction.  The same is true on every central carrier of a finite
product of corner forms.  The native product is a finite weighted
polyhedral intersection with the coefficient restricted to each such
carrier; radial tangency remains true also on its lower-dimensional
faces.  Contracting its tangential coefficient gives exactly the
graded finite rule (CE2).  This verifies the rules for every term of
the supplied polynomial and every common refinement.

If two supplied expressions represent the same native actual form,
their polyhedral currents agree on a common refinement.  Tangential
contraction of their difference is zero there.  The just-verified
finite expression for that contraction is already native actual, so
it represents zero.  This use of currents only checks an equality of
explicit actual expressions; it never obtains an actual representative
from an unspecified current.  It also shows that the operation is
unchanged by further central refinement.
\end{proof}

Let $N$ be the even derivation which sends $d'h_i$ to $d''h_i$ and
kills $h_i,d''h_i,Dh_i$; let $N^+$ be the opposite-colour derivation.
These are the same native exterior operators as frozen HW1.
Define the commuting Euler-weight derivations $E',E''$ by the table
\[
\begin{array}{c|cccc}
 &h_i&d'h_i&d''h_i&Dh_i\\\hline
 E'&h_i&d'h_i&0&0\\
 E''&h_i&0&d''h_i&0.
\end{array}
\tag{CE3}
\]
The full differential and contraction rules give
\[
\begin{gathered}
 d'I'+I'd'=E',\qquad d''I''+I''d''=E'',\\
 d''I'+I'd''=N,\qquad d'I''+I''d'=N^+,\\
 I'^2=I''{}^2=0,\qquad [N,d']=d'',\qquad[N,d'']=0.
\end{gathered}
\tag{CE4}
\]
Indeed each identity holds on each of the four generators in (CE1),
and its two sides are graded derivations of the indicated degree.
The identities consequently hold for every supplied polynomial,
including its full corner currents.  The numerical weights count
scalar factors and the indicated first-colour factors, not the total
bidegree.  A corner factor has weight $(0,0)$ in (CE3).

\subsection{An all-codimension finite solver}

Fix $\ell\ge0$.  Let $u=(A,B,M)$ be an actual triple in
$L_2^{\ell+2}$, so its bidegrees are
\[
 (\ell,\ell+2),\qquad(\ell+2,\ell),\qquad
                         (\ell+1,\ell+1).
 \tag{CE5}
\]
Assume the following finite-presentation conditions:
\begin{enumerate}
\item $A$ and $B$ have pure carrier codimension $\ell$ and
      tangential coefficient degrees $(0,2)$ and $(2,0)$ respectively.
\item $M$ has only carrier codimensions $\ell$ and $\ell+1$,
      with tangential coefficient degrees $(1,1)$ and $(0,0)$ respectively.
\item The Euler weights of $(A,B,M)$ are respectively
      $(1,3),(3,1),(2,2)$.
\item The whole actual triple is closed, and its two additional curls vanish:
\[
 d'A+d''M=0,\qquad d''B+d'M=0,\qquad
                         d''A=0,\quad d'B=0.
 \tag{CE6}
\]
\end{enumerate}
All these are hypotheses about the entire actual expressions, not
chamberwise or probe equalities.  In particular the last two equations
are not inferred merely from the first two.

\begin{theorem}[Explicit two-colour Euler solver]
Set
\[
 P_0=\frac13I''A,\qquad Q_0=\frac13I'B,\qquad
 M_0=M-d'P_0-d''Q_0,\qquad P_1=\frac12I'M_0.
 \tag{CE7}
\]
Every expression is actual and finite, and
\[
 u=\partial_3(0,0,P_0+P_1,Q_0),\qquad
 \partial_3(U,V,P,Q)=(d'U+d''P,d''V+d'Q,d'P+d''Q).
 \tag{CE8}
\]
The denominators are the same in every $\ell$.
\end{theorem}

\begin{proof}
The weights and the last two hypotheses in (CE6) imply by (CE4)
\[
                         d''P_0=A,\qquad d'Q_0=B.
 \tag{CE9}
\]
Use the first two equations of (CE6), anticommutation, and (CE9):
\[
 d''M_0=d''M+d'd''P_0=0,\qquad
 d'M_0=d'M+d''d'Q_0=0.
 \tag{CE10}
\]
The weights of $M_0$ remain $(2,2)$.  The differentiated terms in
(CE7) have only the same two carrier codimensions as $M$: a first
derivative of the tangential coefficient may increase the carrier
codimension by one, and no negative tangential degree can occur.
This also follows directly by differentiating the finite monomials
in (CE1).

Now $NM_0$ has pure carrier codimension $\ell$ with tangential
degree $(0,2)$; $N$ kills the scalar coefficient on the
codimension-$(\ell+1)$ part.  Hence
\[
 I'NM_0=0,
 \qquad d'NM_0=Nd'M_0-d''M_0=0.
 \tag{CE11}
\]
Its prime Euler weight is one, since $N$ lowers that weight by
one.  Applying the first identity of (CE4) to (CE11) gives
$NM_0=0$.  Applying (CE4) again, now to $M_0$, yields
\[
                d'P_1=M_0,\qquad d''P_1=\tfrac12NM_0=0.
 \tag{CE12}
\]
Equations (CE9)--(CE12) are exactly the three components of (CE8).
The coefficients $1/3$ and $1/2$ divide the displayed positive
Euler weights; there is no division by a curvature form or an
operator inverse chosen in a quotient.
\end{proof}

\subsection{The faithful rank-zero-carrier curl test}

Here is the elementary fact that supplies the curl hypotheses in
the rank-three calibration, but not automatically in higher rank.

\begin{lemma}[Prime-closed native forms with no prime degree]
Let $\Omega$ be connected and let $\omega$ be a native actual form
of bidegree $(0,k)$ on the faithful full coordinate chart $\Omega$.
If $d'\omega=0$, all its coefficients are constant.  The analogous
assertion holds after exchanging the two colours.
\end{lemma}

\begin{proof}
There can be no positive carrier codimension in degree $(0,k)$.
Thus $\omega=\sum_J u_J(x)d''x_J$ is a piecewise-smooth,
locally integrable coefficient form.  In the faithful full chart,
$d'\omega=0$ says that every distributional derivative
$\partial u_J/\partial x_i$ is zero.  This includes the jumps on
every wall, not only derivatives in chamber interiors.  To see
directly that $u_J$ is constant, convolve on a small ball with a
smooth compactly supported mollifier.  Every derivative of the
convolution is zero on the smaller ball, so the convolution is
constant there.  Its distributional limit is a constant on that
ball; overlapping balls identify these constants, and connectedness
gives the assertion.  The resulting identity is an equality of the
already native actual coefficient forms.
\end{proof}

For $\ell=0$, a closed triple satisfies
$d'd''A=0$ and $d''d'B=0$.  The lemma makes the two curls
constant.  Their constants must still be determined on a genuine
ordinary open with all relevant coordinate directions free; they
need not be zero.  The cubic unit triple
\[
 E_{abc}=(-a\,d''b\wedge d''c,\ b\,d'a\wedge d'c,
                                      c\,d'a\wedge d''b)
 \tag{CE13}
\]
has nonzero constant first curl
$-d''a\wedge d''b\wedge d''c$.  Subtracting that extending
triple is essential in the raw-cube application.

For $\ell>0$, the first curl has positive carrier codimension and
the constant-coefficient lemma does not apply.  It may be a nonzero
closed polyhedral current even with the required homogeneous weight.
In particular (CE8) cannot be invoked for the full higher-rank cube
without separately proving both curl identities in (CE6).  The
rank-four retained-GR-remainder calibration already warns against
assuming these identities from the rank-three endpoint.

\paragraph{Exact actual interface and limits.}
The native algebra, equivalence, local carrier description, refinement
compatibility, and finite operations have been specified above.
No new completed or dual form space is used.  On one fixed finite
presentation, (CE2)--(CE8) are finite coefficient operations with the
usual finite mixed Leibniz bounds.  There is no assertion of a
continuous contraction on the whole actual space or a whole-kernel
topology theorem.  To use the result for Berkovich data one must
rewrite (CE2) entirely as finite scalar/differential operations in a
supplied native presentation and verify substitution compatibility.
The separate NCP positive-unit-log and norm-preservation lemmas give
one precise such interface for the capped norm expressions; they do
not by themselves supply the missing higher-rank curl equations.

The conical condition is substantive.  For a cap with a fixed nonzero
real offset, its walls may be affine rather than central, and radial
contraction of its first derivative need not be an actual scalar.
No such contraction is licensed here.  Likewise no current is
substituted for the actual Green form required on the original
support complement.
